%% Dynamic LoR-GS Project Report
%% CSC 2210 - Visual and Mobile Computing, University of Toronto
%%
\documentclass[sigconf]{acmart}

\AtBeginDocument{%
  \providecommand\BibTeX{{Bib\TeX}}}

\setcopyright{none}
\acmConference[CSC 2210]{Visual and Mobile Computing}{2025}{University of Toronto}
\settopmatter{printacmref=false}
\renewcommand\footnotetextcopyrightpermission[1]{}
\pagestyle{plain}

\usepackage{graphicx}
\usepackage{booktabs}
\usepackage{amsmath}
\usepackage{amssymb}
\usepackage{microtype}
\usepackage{xcolor}


\newcommand{\TODO}[1]{\textcolor{red}{\textbf{[#1]}}}

%% -------------------------------------------------------
\begin{document}

\title{Dynamic LoR-GS: Adaptive Low-Rank Gaussian Parameters\\
       for Memory-Efficient 3D Gaussian Splatting Training}

\author{\TODO{Author One}}
\affiliation{\institution{University of Toronto}}
\email{author1@mail.utoronto.ca}

\author{\TODO{Author Two}}
\affiliation{\institution{University of Toronto}}
\email{author2@mail.utoronto.ca}

\author{\TODO{Author Three}}
\affiliation{\institution{University of Toronto}}
\email{author3@mail.utoronto.ca}

\begin{abstract}
\TODO{Fill in last}
\end{abstract}


\maketitle

%% -------------------------------------------------------
\section{Background and Related Work}
\label{sec:background}

3D Gaussian Splatting (3DGS)~\cite{kerbl2023} has rapidly become the dominant
paradigm for novel view synthesis, offering high visual quality with a real-time rendering 
and fast training. By representing scenes as collections of explicit, differentiable 
3D Gaussian primitives rendered via fast tile-based rasterization, 
3DGS achieves high visual fidelity with training times measured in minutes. 
This has sparked a large body of follow-on work spanning
dynamic scenes~\cite{lr4dgstream}, scene editing~\cite{gaussianeditor}, compression~\cite{Papantonakis_2024, survey3dgs}, 
and avatar generation~\cite{avatar}.

Despite this success, training 3DGS on memory-constrained hardware remains
difficult~\cite{kerbl2023, Papantonakis_2024, survey3dgs}. Peak GPU memory during 
training is driven by two coupled factors.
First, 3DGS employs an adaptive densification strategy that continuously clones
and splits Gaussian primitives throughout optimization, growing the primitive count
$N$ from an initial sparse point cloud to several million --- each new Gaussian
adding to every active parameter tensor. Second, each Gaussian carries a substantial
per-primitive parameter budget. Beyond position, opacity, and covariance, the
view-dependent color of each primitive is encoded by spherical harmonic (SH)
coefficients. Using degree-3 SH, each Gaussian stores $(3+1)^2 = 16$ coefficients
per RGB channel, totaling 48 floating-point values per Gaussian. As noted in prior
memory analyses~\cite{Papantonakis_2024, survey3dgs}, these SH coefficients alone account
for more than half of all model parameters --- for a scene with $N = 3$ million
Gaussians, this amounts to roughly 576\,MB before optimizer states and gradients
are accounted for. The combination of a growing $N$ and a large per-Gaussian
footprint makes peak training VRAM a hard constraint that existing methods have
largely failed to address.

A natural response is to reduce $N$ directly. Pruning-based methods such as
Mini-Splatting~\cite{minisplatting} and LP-3DGS~\cite{lp3dgs} remove redundant
Gaussians to produce more compact representations. However, aggressive pruning often
trades reconstruction quality for memory, since each Gaussian encodes fine
geometric and appearance detail that cannot always be recovered once removed.
We take a different position: rather than reducing the number of primitives,
we ask whether the \emph{per-Gaussian parameter cost} can be reduced adaptively
--- allocating more capacity to primitives that genuinely require it and less
to those that do not --- without compromising the Gaussian count or the
reconstruction quality it affords.

This question has been explored in prior work, but only after training completes.
Reduced3DGS~\cite{Papantonakis_2024} introduces an adaptive per-Gaussian SH band
selection based on view-dependent transmittance variance, substantially reducing
storage and rendering memory. Critically, the authors acknowledge that this
approach does not reduce \emph{peak training memory}, since the full SH tensor
must be maintained throughout optimization before any pruning can occur.
LFGS~\cite{lfgs} similarly applies dynamic SH masking post-training for
compression. SG-Splatting~\cite{sgsplatting} and MEGS2~\cite{megs2} replace
SH with spherical Gaussian lobes for more compact rendering, but again leave
the training regime unchanged. The gap is clear: no existing method reduces
per-Gaussian parameter cost \emph{during} training, where peak VRAM is
actually determined.

Our approach is to apply Low-Rank Adaptation (LoRA)~\cite{lora} directly to
the learnable parameters of each Gaussian primitive during training. LoRA,
originally proposed for parameter-efficient fine-tuning of large language models,
replaces a parameter matrix $W \in \mathbb{R}^{m \times n}$ with a low-rank
factorization $W \approx AB$, where $A \in \mathbb{R}^{m \times r}$,
$B \in \mathbb{R}^{r \times n}$, and $r \ll \min(m, n)$, reducing parameter
count from $mn$ to $r(m + n)$. By maintaining only the low-rank factors
throughout training, the full parameter matrix is never allocated, directly
reducing peak VRAM rather than compressing after the fact.

Among the Gaussian attributes, the SH coefficient tensor is the natural primary
candidate for low-rank factorization. Its structure as a matrix
$C \in \mathbb{R}^{3 \times 16}$ (RGB channels $\times$ SH basis functions)
makes the decomposition well-defined, and its dominance in the per-Gaussian
memory budget makes it the highest-leverage target. The observation that
Reduced3DGS~\cite{Papantonakis_2024} can prune high-degree SH bands from a large
fraction of Gaussians with minimal quality loss supports our hypothesis that
the effective rank of $C$ is low for most primitives. Furthermore, 
rather than applying a fixed rank across all primitives, we draw inspiration from adaptive 
budget allocation in language models~\cite{adalora} to dynamically distribute rank. 
A flat matte surface requires far less representational capacity than a complex specular 
reflection; our method dynamically models this heterogeneity.
We also explore LoRA applied to other Gaussian attributes to characterize where low-rank structure
is and is not beneficial; these experiments and their outcomes are reported in
Section~\ref{sec:method}.

LoRA has seen use in the 3DGS ecosystem, but not for this purpose.
LR-4DGStream~\cite{lr4dgstream} applies low-rank adaptation to a deformation
model for dynamic scene streaming bandwidth reduction. UFV-Splatter~\cite{ufvsplatter}
injects LoRA into a feed-forward Gaussian-prediction network to handle unfavorable
camera poses. In both cases, LoRA is applied to auxiliary neural network weights,
not to Gaussian primitive parameters directly. StyleSplat~\cite{stylesplat} comes
closest in spirit, directly optimizing the SH coefficients of selected Gaussians
via a feature-matching loss to transfer artistic style --- but without any
low-rank constraint or memory reduction goal. To our knowledge, no prior work
applies low-rank factorization to the parameters of Gaussian primitives during
training as a mechanism for reducing peak GPU memory.

%% -------------------------------------------------------
\section{Goal}
\label{sec:goal}

he primary goal of this project is to drastically reduce peak GPU training memory 
in 3DGS, enabling high-fidelity scene optimization on memory-constrained, consumer-grade 
hardware. We achieve this by replacing full-rank Gaussian parameters with per-Gaussian 
low-rank factorizations. By dynamically allocating rank across primitives based on 
their estimated view-dependent complexity, we aim to eliminate standard memory 
bottlenecks without reducing the overall Gaussian count or sacrificing rendering quality.

%% -------------------------------------------------------
\section{Methodology}
\label{sec:method}

\TODO{Fill in}

%% -------------------------------------------------------
\section{Observations and Motivation}
\label{sec:observations}

\TODO{Fill in}

%% -------------------------------------------------------
\section{Ideas and Solutions}
\label{sec:solutions}

\TODO{Fill in}

%% -------------------------------------------------------
\section{Results}
\label{sec:results}

\TODO{Fill in}

%% -------------------------------------------------------
\section{Takeaways}
\label{sec:takeaways}

\TODO{Fill in}

%% -------------------------------------------------------
\bibliographystyle{ACM-Reference-Format}
\bibliography{references}

\end{document}